\section{Related Work}\label{sec:related}
Prior work on RPL security, mainly from 2018--2026 with earlier foundational sources for RPL and the methods used here, falls into four strands: DIS flooding and its defenses, DAO attacks and combined attacks, adaptive and cooperative detection, and learning-based RPL security. Table~\ref{tab:positioning} contrasts the closest methods with this study, and Section~\ref{sec:gaps} states the gaps they leave open.

\subsection{RPL Control-Plane Attacks and DIS Flooding}
RFC~6550 defines the DIS, DIO, and DAO messages and leaves their authentication optional \cite{rfc6550}. Raoof et al.\ survey the resulting attacks and their mitigations \cite{raoof2019}, and Albinali and Azzedin classify them by what the attacker does with a control message: generating, dropping, modifying, or replaying it \cite{albinali2024}. Many of these attacks have been studied on their own, among them Sybil \cite{khan2023}, black hole \cite{sharma2022}, and flooding attacks in cloud-connected low-power networks \cite{ankam2023}.

Rajasekar and Rajkumar measured the impact of DIS flooding, the most studied resource attack, on 6LoWPAN network performance \cite{rajasekar2022}. Secure-RPL limits DIS processing by rejecting solicitations that arrive too early \cite{vermaranga2020}; Abhinaya and Sudhakar proposed a secure RPL variant against DIS flooding \cite{abhinaya2021}; and RPL-MRC delays DIO responses to multicast DIS \cite{medjek2021}. DIS-FDM, our base method, moved from thresholds to evidence: each node fuses a DIS traffic-load indicator and a parent-change indicator with Dempster's rule, adds an energy indicator when the result is inconclusive, and blacklists a neighbor through counter rules backed by the fused belief \cite{prajisha2025}. Its DIS-load threshold uses the Beta-mean estimate of Nishanth and Mujeeb, who modeled flooding attacks in ad hoc networks \cite{nishanth2021}. DIS-FDM reduced DIO overhead in Contiki/Cooja, but its belief masses are lookup values fixed after empirical tuning, its published results report no detection rates, and its evaluation does not include repeated-run statistics \cite{prajisha2025}. All of these defenses observe only DIS and DIO traffic.

\subsection{DAO Attacks and Combined Attacks}
On the downward routes, Pu studied the DAO inconsistency attack and proposed a lightweight mitigation \cite{pu2018}. Ghaleb et al.\ \cite{ghaleb2019} and Li-MSD \cite{limsd2024} addressed the DAO insider attack, in which a node sends excessive DAOs; Li-MSD blacklists a node whose DAO count exceeds a threshold. Rouissat et al.\ studied DIO inundation \cite{rouissat2023}. IMDRPL protects DIO and DAO exchanges with certificate attestation \cite{imdrpl2023}, which keeps outsiders away but requires an identity infrastructure and does not stop an insider that holds valid credentials. Counting defenses such as Li-MSD detect excessive DAO volume, not a well-formed DAO that names the wrong parent, which is the manipulation we study.

Alsukayti and Alreshoodi measured RPL under simple and complex (combined) routing attacks \cite{alsukayti2023}. DE2RA-RPL handles DAO flooding together with rank and version attacks \cite{de2rarpl2024}, and Abdulkareem et al.\ detected joint rank, wormhole, and version attacks with a collaborative IDS \cite{abdulkareem2024}. Garcia Ribera et al.\ tested a hybrid IDS against several RPL attacks, including flooding attacks run concurrently by different nodes \cite{garciaribera2022}. In all of these studies, either the combined vectors differ from DIS and DAO, or each attacker performs one attack. None evaluates a node that floods DIS and manipulates its DAO parent at once.

\subsection{Adaptive, Cooperative, and Hybrid Detection}
FLSec-RPL adapts its decision with fuzzy logic against DIO neighbor suppression in static and mobile networks \cite{flsecrpl2024}. DSAT-IDS derives adaptive thresholds from subtree forwarding behavior against selective forwarding \cite{dsatids2024}. Cooperative and hybrid IDSs share observations between nodes or with a central component against rank and other attacks \cite{monteiro2026,garciaribera2022}, and Azzaoui et al.\ run lightweight neural detectors on RPL parents that report to a central system \cite{azzaoui2024}. These systems show the value of adaptive thresholds and shared evidence, but their observations (forwarding rates, rank values, DIO suppression) do not reveal a false DAO parent. Reputation systems offer a principled alternative to fixed beliefs. The Beta reputation of Jøsang and Ismail \cite{josang2002}, which Ganeriwal et al.\ adapted to sensor networks \cite{ganeriwal2008}, estimates trust as the posterior mean of a Beta distribution updated from observed behavior. We use this idea to calibrate the DAO evidence.

\subsection{Learning-Based RPL Security}
QSec-RPL applies Q-learning to detect version-number attacks in mobile RPL \cite{qsecrpl2023}. iTRPL combines trust with multi-agent reinforcement learning, dividing work between nodes and the root \cite{itrpl2024}. CRRDS runs a deep reinforcement learning agent at the root against decreased-rank attacks \cite{haqdad2026}, and Ahmadi and Javidan proposed a distributed mitigation for multiple routing attacks that trains its models by federated learning \cite{ahmadi2026}. Panda and Muthuraman trained supervised classifiers on a simulation-driven RPL attack dataset \cite{panda2026}. Outside RPL, Dempster--Shafer evidence has been paired with reinforcement learning in vehicular reputation systems coordinated by edge servers or roadside units \cite{gyawali2021,jasti2024}. These works target other attacks, learn at the root, at a server, or through federated training, or classify offline. None places a learner on each constrained node over fused DIS and DAO evidence.

\begin{table*}[t]
\caption{Positioning of AEF-RPL relative to the closest RPL security methods.}\label{tab:positioning}
\centering\footnotesize
\setlength{\tabcolsep}{4pt}
\renewcommand{\arraystretch}{1.12}
\begin{tabularx}{\textwidth}{@{}p{0.13\textwidth}p{0.15\textwidth}p{0.19\textwidth}p{0.17\textwidth}X@{}}
\toprule
Method & Threat & Mechanism & Limitation for the combined threat & How this study differs \\
\midrule
DIS-FDM \cite{prajisha2025} & Multicast DIS flooding & Local Dempster--Shafer fusion of DIS, parent-change, and energy evidence & Blind to DAO; fixed belief masses; no accuracy or statistics reported & Adds self-calibrated DAO evidence and repair; reports TPR, FPR, and paired tests \\
Secure-RPL \cite{vermaranga2020} & DIS flooding & Rejects early DIS arrivals & No DAO view; no learning & Evaluated as a baseline under the combined attack \\
RPL-MRC \cite{medjek2021} & Multicast DIS & Delays DIO responses & No DAO view; no learning & Evaluated as a baseline \\
Li-MSD \cite{limsd2024} & DAO insider (volume) & DAO counting and blacklisting & Misses a well-formed DAO with a false parent & Checks DAO content (parent) instead of volume \\
Garcia Ribera et al.\ \cite{garciaribera2022} & Several attacks, concurrent flooding & Hybrid distributed and central IDS & One attack per node; central component & Same node runs both attacks; fully local decisions \\
DE2RA-RPL \cite{de2rarpl2024}, Abdulkareem et al.\ \cite{abdulkareem2024} & DAO flood with rank/version; joint rank, wormhole, version & Combined detection rules; collaborative IDS & Different vector combinations & Targets the DIS plus DAO-parent combination \\
QSec-RPL \cite{qsecrpl2023}, iTRPL \cite{itrpl2024} & Version number; insider attacks & Q-learning; trust with multi-agent RL & Other attacks; root-side decision (iTRPL) & Per-node tabular learner over fused DIS and DAO evidence \\
CRRDS \cite{haqdad2026} & Decreased rank & Deep RL at the root & Root resources and central view & 780-byte learner on every constrained node \\
\midrule
AEF-RPL (this study) & Same-node DIS flood and false-parent DAO & Fused DIS, parent, energy, and calibrated DAO evidence; selective repair; on-node Q-learning & Evaluated in 30-node simulation & Paired comparison against all of the above, with learning isolated by ablation \\
\bottomrule
\end{tabularx}
\end{table*}

\subsection{Research Gaps}\label{sec:gaps}
To our knowledge, no RPL defense has been tested against an adversary that runs multicast DIS flooding and DAO manipulation at the same time from the same node (Gap~1). Defenses for the DIS plane and for the DAO plane have developed as separate lines of work \cite{vermaranga2020,medjek2021,prajisha2025,limsd2024}, IDSs that handle several attacks place each attack on a different node \cite{garciaribera2022,imdrpl2023}, and the combined-attack studies we found pair other vectors \cite{alsukayti2023,de2rarpl2024,abdulkareem2024}. The DIS-FDM authors themselves leave this hybrid scenario to future work \cite{prajisha2025}.

The belief masses in DIS-FDM are fixed after offline empirical tuning and do not change with how reliable an observation turns out to be in a given network (Gap~2).

DIS-FDM's method section announces a detection-accuracy measurement, yet the published results give no true- or false-positive rates and no statistical tests over repeated runs \cite{prajisha2025} (Gap~3). Without these, a gain in delivery cannot be told apart from simulation noise or from the isolation of honest nodes.

The lightweight DIS and DAO defenses rely on fixed rules and thresholds, so a node cannot learn from its own experience how to treat a suspicious neighbor (Gap~4). The existing learning-based RPL defenses address other attacks and learn at the root \cite{itrpl2024,haqdad2026}, train models offline or by federation \cite{ahmadi2026,panda2026}, or do not use DAO evidence \cite{qsecrpl2023,azzaoui2024}, and Dempster--Shafer evidence has been combined with reinforcement learning only in vehicular systems coordinated by edge servers or roadside units \cite{gyawali2021,jasti2024}.
